\subsection{\label{sec.gf.weighted-set}The generating function of a weighted
set}

So far, we have built a theory of FPSs, and we have occasionally applied them
to counting problems. However, so far, the latter applications were either
limited to counting tuples of integers (because such tuples naturally appear
as indexes when we multiply several sums) or somewhat indirect, mediated by
integer sequences (i.e., we could not immediately transform our counting
problem into an FPS but rather had to first define an integer sequence and
then take its generating function). For instance, our above proof of Theorem
\ref{thm.gf.prod.euler-comb} was an application of the former type, while
Example 2 in Section \ref{sec.gf.exas} was one of the latter type.

I shall now explain a more direct approach to transforming counting problems
into FPSs. This approach is the theory of so-called \textquotedblleft%
\emph{combinatorial classes}\textquotedblright\ or \textquotedblleft%
\emph{finite-type weighted sets}\textquotedblright, and can be used
(particularly in some of its more advanced variants) to prove deep results.
However, I will not use it much in these notes, so I will only give a brief
introduction. Much more can be learned from \cite[Part A]{FlaSed09};
introductions can also be found in \cite[\S 3.3--\S 3.4]{Fink17}, \cite[Parts
1 and 2]{Melcze24} and \cite[Chapter 6]{Loehr-BC}. The theory is not magic,
and in principle can always be eschewed (any proof using it can be rewritten
without its use), but it offers the major advantages of systematizing the
combinatorial objects being counted and making the proofs more insightful. In
a sense, it allows us to perform certain computations not just with the
generating functions of the objects, but with the objects themselves -- a
middle ground between algebra and combinatorics.

I will now outline the beginnings of this theory, following Fink's
introduction \cite[\S 3.3--\S 3.4]{Fink17} (but speaking of \emph{finite-type
weighted sets} where \cite{Fink17}, \cite{Melcze24} and \cite{FlaSed09} speak
of \textquotedblleft combinatorial classes\textquotedblright).

\subsubsection{The theory}

\begin{definition}
\label{def.gf-ws.weighted-sets}\textbf{(a)} A \emph{weighted set} is a set $A$
equipped with a function $w:A\rightarrow\mathbb{N}$, which is called the
\emph{weight function} of this weighted set. For each $a\in A$, the value
$w\left(  a\right)  $ is denoted $\left\vert a\right\vert $ and is called the
\emph{weight} of $a$ (in our weighted set). \medskip

Usually, instead of explicitly specifying the weight function $w$ as a
function, we will simply specify the weight $\left\vert a\right\vert $ for
each $a\in A$. The weighted set consisting of the set $A$ and the weight
function $w$ will be called $\left(  A,w\right)  $ or simply $A$ when the
weight function $w$ is clear from the context. \medskip

\textbf{(b)} A weighted set $A$ is said to be \emph{finite-type} if for each
$n\in\mathbb{N}$, there are only finitely many $a\in A$ having weight
$\left\vert a\right\vert =n$. \medskip

\textbf{(c)} If $A$ is a finite-type weighted set, then the \emph{weight
generating function} of $A$ is defined to be the FPS%
\[
\sum_{a\in A}x^{\left\vert a\right\vert }=\sum_{n\in\mathbb{N}}\left(
\text{\# of }a\in A\text{ having weight }n\right)  \cdot x^{n}\in
\mathbb{Z}\left[  \left[  x\right]  \right]  .
\]
This FPS is denoted by $\overline{A}$. \medskip

\textbf{(d)} An \emph{isomorphism} between two weighted sets $A$ and $B$ means
a bijection $\rho:A\rightarrow B$ that preserves the weight (i.e., each $a\in
A$ satisfies $\left\vert \rho\left(  a\right)  \right\vert =\left\vert
a\right\vert $). \medskip

\textbf{(e)} We say that two weighted sets $A$ and $B$ are \emph{isomorphic}
(this is written $A\cong B$) if there exists an isomorphism between $A$ and
$B$.
\end{definition}

\begin{example}
\label{exa.ws.bin-string1}Let $B$ be the weighted set of all \emph{binary
strings}, i.e., finite tuples consisting of $0$s and $1$s. Thus,%
\[
B=\left\{  \left(  {}\right)  ,\ \left(  0\right)  ,\ \left(  1\right)
,\ \left(  0,0\right)  ,\ \left(  0,1\right)  ,\ \left(  1,0\right)
,\ \left(  1,1\right)  ,\ \left(  0,0,0\right)  ,\ \left(  0,0,1\right)
,\ \ldots\right\}  .
\]
The weight of an $n$-tuple is defined to be $n$. In other words, the weight of
a binary string is the length of this string. This weighted set $B$ is
finite-type, since for each $n\in\mathbb{N}$, there are only finitely many
binary strings of length $n$ (namely, there are $2^{n}$ such strings). The
weight generating function of $B$ is%
\[
\overline{B}=\sum_{n\in\mathbb{N}}\underbrace{\left(  \text{\# of }a\in
B\text{ having weight }n\right)  }_{\substack{=\left(  \text{\# of binary
strings of length }n\right)  \\=2^{n}}}\cdot\,x^{n}=\sum_{n\in\mathbb{N}}%
2^{n}x^{n}=\dfrac{1}{1-2x}.
\]

\end{example}

\begin{example}
Let $B^{\prime}$ be the weighted set of all binary strings, i.e., finite
tuples consisting of $0$s and $1$s. The weight of an $n$-tuple $\left(
a_{1},a_{2},\ldots,a_{n}\right)  \in B^{\prime}$ is defined to be $a_{1}%
+a_{2}+\cdots+a_{n}$. (Thus, $B^{\prime}$ is the same set as the $B$ in
Example \ref{exa.ws.bin-string1}, but equipped with a different weight
function.) This weighted set $B^{\prime}$ is \textbf{not} finite-type, since
there are infinitely many $a\in B^{\prime}$ having weight $0$ (namely, all
tuples of the form $\left(  0,0,\ldots,0\right)  $, with any number of zeroes.)
\end{example}

Those familiar with graded vector spaces (or graded modules) will recognize
weighted sets as their combinatorial analogues: Weighted sets are to sets what
graded vector spaces are to vector spaces. (In particular, the weight
generating function $\overline{A}$ of a finite-type weighted set $A$ is an
analogue of the Hilbert series of a finite-type graded vector space.) \medskip

Let us now prove a few basic properties of weighted sets and their weight
generating functions.

\begin{proposition}
\label{prop.gf-ws.iso}Let $A$ and $B$ be two isomorphic finite-type weighted
sets. Then, $\overline{A}=\overline{B}$.
\end{proposition}

\begin{proof}
This is almost trivial: The weighted sets $A$ and $B$ are isomorphic; thus,
there exists an isomorphism $\rho:A\rightarrow B$. Consider this $\rho$. Then,
$\rho$ is a bijection and preserves the weight (since $\rho$ is an isomorphism
of weighted sets). The latter property says that we have
\begin{equation}
\left\vert \rho\left(  a\right)  \right\vert =\left\vert a\right\vert
\ \ \ \ \ \ \ \ \ \ \text{for each }a\in A. \label{pf.prop.gf-ws.iso.wtpres}%
\end{equation}
Now, the definition of $\overline{B}$ yields%
\begin{align*}
\overline{B}  &  =\sum_{b\in B}x^{\left\vert b\right\vert }=\sum_{a\in
A}\underbrace{x^{\left\vert \rho\left(  a\right)  \right\vert }}%
_{\substack{=x^{\left\vert a\right\vert }\\\text{(by
(\ref{pf.prop.gf-ws.iso.wtpres}))}}}\ \ \ \ \ \ \ \ \ \ \left(
\begin{array}
[c]{c}%
\text{here, we have substituted }\rho\left(  a\right)  \text{ for }b\\
\text{in the sum, since }\rho\text{ is a bijection}%
\end{array}
\right) \\
&  =\sum_{a\in A}x^{\left\vert a\right\vert }.
\end{align*}
Comparing this with%
\[
\overline{A}=\sum_{a\in A}x^{\left\vert a\right\vert }%
\ \ \ \ \ \ \ \ \ \ \left(  \text{by the definition of }\overline{A}\right)
,
\]
we obtain $\overline{A}=\overline{B}$. This proves Proposition
\ref{prop.gf-ws.iso}.
\end{proof}

Note that Proposition \ref{prop.gf-ws.iso} has a converse: If $\overline
{A}=\overline{B}$, then $A\cong B$. \medskip

Recall that the \emph{disjoint union} of two sets $A$ and $B$ is informally
understood to be \textquotedblleft the union of $A$ and $B$ where we pretend
that the sets $A$ and $B$ are disjoint (even if they aren't)\textquotedblright%
. Formally, it is defined as the set $\left(  \left\{  0\right\}  \times
A\right)  \cup\left(  \left\{  1\right\}  \times B\right)  $, but we think of
the elements $\left(  0,a\right)  $ of this set as copies of the respective
elements $a\in A$, and we think of the elements $\left(  1,b\right)  $ of this
set as copies of the respective elements $b\in B$. (Thus, even if the original
sets $A$ and $B$ have some elements in common, say $c\in A\cap B$, the two
copies $\left(  0,c\right)  $ and $\left(  1,c\right)  $ of such common
elements will not coincide.)

If $A$ and $B$ are two finite sets, then their disjoint union always has size
$\left\vert A\right\vert +\left\vert B\right\vert $.

Let us now extend the concept of a disjoint union to weighted sets:

\begin{definition}
\label{def.gf-ws.djun}Let $A$ and $B$ be two weighted sets. Then, the weighted
set $A+B$ is defined to be the disjoint union of $A$ and $B$, with the weight
function inherited from $A$ and $B$ (meaning that each element of $A$ has the
same weight that it had in $A$, and each element of $B$ has the same weight
that it had in $B$). Formally speaking, this means that $A+B$ is the set
$\left(  \left\{  0\right\}  \times A\right)  \cup\left(  \left\{  1\right\}
\times B\right)  $, with the weight function given by
\begin{equation}
\left\vert \left(  0,a\right)  \right\vert =\left\vert a\right\vert
\ \ \ \ \ \ \ \ \ \ \text{for each }a\in A \label{eq.def.gf-ws.djun.wt0}%
\end{equation}
and%
\begin{equation}
\left\vert \left(  1,b\right)  \right\vert =\left\vert b\right\vert
\ \ \ \ \ \ \ \ \ \ \text{for each }b\in B. \label{eq.def.gf-ws.djun.wt1}%
\end{equation}

\end{definition}

\begin{proposition}
\label{prop.gf-ws.djun}Let $A$ and $B$ be two finite-type weighted sets. Then,
$A+B$ is finite-type, too, and satisfies $\overline{A+B}=\overline
{A}+\overline{B}$.
\end{proposition}

\begin{proof}
The formal definition of $A+B$ says that $A+B=\left(  \left\{  0\right\}
\times A\right)  \cup\left(  \left\{  1\right\}  \times B\right)  $. Thus, the
set $A+B$ is the union of the two disjoint sets $\left\{  0\right\}  \times A$
and $\left\{  1\right\}  \times B$ (indeed, these two sets are clearly
disjoint, since $0\neq1$).

Let us first check that $A+B$ is finite-type.

For each $n\in\mathbb{N}$, there are only finitely many $a\in A$ having weight
$\left\vert a\right\vert =n$ (since $A$ is finite-type), and there are only
finitely many $b\in B$ having weight $\left\vert b\right\vert =n$ (since $B$
is finite-type). Hence, there are only finitely many $c\in A+B$ having weight
$\left\vert c\right\vert =n$ (because any such $c$ either has the form
$\left(  0,a\right)  $ for some $a\in A$ having weight $\left\vert
a\right\vert =n$, or has the form $\left(  1,b\right)  $ for some $b\in B$
having weight $\left\vert b\right\vert =n$). In other words, the weighted set
$A+B$ is finite-type.

Let us now check that $\overline{A+B}=\overline{A}+\overline{B}$. Indeed, the
definition of $\overline{A+B}$ yields%
\begin{align*}
\overline{A+B}  &  =\sum_{c\in A+B}x^{\left\vert c\right\vert }\\
&  =\underbrace{\sum_{c\in\left\{  0\right\}  \times A}x^{\left\vert
c\right\vert }}_{\substack{=\sum_{a\in A}x^{\left\vert \left(  0,a\right)
\right\vert }\\\text{(here, we have substituted }\left(  0,a\right)
\\\text{for }c\text{ in the sum, since the}\\\text{map }A\rightarrow\left\{
0\right\}  \times A,\ a\mapsto\left(  0,a\right)  \text{ is a bijection)}%
}}+\underbrace{\sum_{c\in\left\{  1\right\}  \times B}x^{\left\vert
c\right\vert }}_{\substack{=\sum_{b\in B}x^{\left\vert \left(  1,b\right)
\right\vert }\\\text{(here, we have substituted }\left(  1,b\right)
\\\text{for }c\text{ in the sum, since the}\\\text{map }B\rightarrow\left\{
1\right\}  \times B,\ b\mapsto\left(  1,b\right)  \text{ is a bijection)}}}\\
&  \ \ \ \ \ \ \ \ \ \ \ \ \ \ \ \ \ \ \ \ \left(
\begin{array}
[c]{c}%
\text{here, we have split the sum, since the}\\
\text{set }A+B\text{ is the union of the two disjoint}\\
\text{sets }\left\{  0\right\}  \times A\text{ and }\left\{  1\right\}  \times
B
\end{array}
\right) \\
&  =\sum_{a\in A}\underbrace{x^{\left\vert \left(  0,a\right)  \right\vert }%
}_{\substack{=x^{\left\vert a\right\vert }\\\text{(by
(\ref{eq.def.gf-ws.djun.wt0}))}}}+\sum_{b\in B}\underbrace{x^{\left\vert
\left(  1,b\right)  \right\vert }}_{\substack{=x^{\left\vert b\right\vert
}\\\text{(by (\ref{eq.def.gf-ws.djun.wt1}))}}}=\underbrace{\sum_{a\in
A}x^{\left\vert a\right\vert }}_{=\overline{A}}+\underbrace{\sum_{b\in
B}x^{\left\vert b\right\vert }}_{=\overline{B}}=\overline{A}+\overline{B}.
\end{align*}
Thus, the proof of Proposition \ref{prop.gf-ws.djun} is complete.
\end{proof}

We can easily extend Definition \ref{def.gf-ws.djun} and Proposition
\ref{prop.gf-ws.djun} to disjoint unions of any number (even infinite) of
weighted sets. (However, in the case of infinite disjoint unions, we need to
require the disjoint union to be finite-type in order for Proposition
\ref{prop.gf-ws.djun} to make sense.)

We note that disjoint unions of weighted sets do not satisfy associativity
\textquotedblleft on the nose\textquotedblright: If $A$, $B$ and $C$ are three
weighted sets, then the weighted sets $A+B+C$ and $\left(  A+B\right)  +C$ and
$A+\left(  B+C\right)  $ are not literally equal but rather are isomorphic via
canonical isomorphisms. But Proposition \ref{prop.gf-ws.iso} shows that their
weight generating functions are the same, so that we need not distinguish
between them if we are only interested in these functions. \medskip

Another operation on weighted sets is the Cartesian product:

\begin{definition}
\label{def.gf-ws.prod}Let $A$ and $B$ be two weighted sets. Then, the weighted
set $A\times B$ is defined to be the Cartesian product of the sets $A$ and $B$
(that is, the set $\left\{  \left(  a,b\right)  \ \mid\ a\in A\text{ and }b\in
B\right\}  $), with the weight function defined as follows: For any $\left(
a,b\right)  \in A\times B$, we set%
\begin{equation}
\left\vert \left(  a,b\right)  \right\vert =\left\vert a\right\vert
+\left\vert b\right\vert . \label{eq.def.gf-ws.prod.wt}%
\end{equation}

\end{definition}

\begin{proposition}
\label{prop.gf-ws.prod}Let $A$ and $B$ be two finite-type weighted sets. Then,
$A\times B$ is finite-type, too, and satisfies $\overline{A\times B}%
=\overline{A}\cdot\overline{B}$.
\end{proposition}

\begin{proof}
[Proof of Proposition \ref{prop.gf-ws.prod} (sketched).]The proof that
$A\times B$ is finite-type is easy\footnote{Here is the idea: Let
$n\in\mathbb{N}$. We must prove that there are only finitely many pairs
$\left(  a,b\right)  \in A\times B$ having weight $\left\vert \left(
a,b\right)  \right\vert =n$. Since $\left\vert \left(  a,b\right)  \right\vert
=\left\vert a\right\vert +\left\vert b\right\vert $, these pairs have the
property that $a\in A$ has weight $k$ for some $k\in\left\{  0,1,\ldots
,n\right\}  $, and that $b\in B$ has weight $n-k$ for the same $k$. This
leaves only finitely many options for $k$, only finitely many options for $a$
(since $A$ is finite-type and thus has only finitely many elements of weight
$k$), and only finitely many options for $b$ (since $B$ is finite-type and
thus has only finitely many elements of $n-k$). Altogether, we thus obtain
only finitely many options for $\left(  a,b\right)  $.}.

Now, the definition of $\overline{A\times B}$ yields%
\begin{align*}
\overline{A\times B}  &  =\sum_{\left(  a,b\right)  \in A\times B}%
\underbrace{x^{\left\vert \left(  a,b\right)  \right\vert }}%
_{\substack{=x^{\left\vert a\right\vert +\left\vert b\right\vert }\\\text{(by
(\ref{eq.def.gf-ws.prod.wt}))}}}\ \ \ \ \ \ \ \ \ \ \left(
\begin{array}
[c]{c}%
\text{since all elements of }A\times B\text{ have}\\
\text{the form }\left(  a,b\right)  \text{ for some }a\text{ and }b
\end{array}
\right) \\
&  =\sum_{\left(  a,b\right)  \in A\times B}\underbrace{x^{\left\vert
a\right\vert +\left\vert b\right\vert }}_{=x^{\left\vert a\right\vert }\cdot
x^{\left\vert b\right\vert }}\\
&  =\sum_{\left(  a,b\right)  \in A\times B}x^{\left\vert a\right\vert }\cdot
x^{\left\vert b\right\vert }=\underbrace{\left(  \sum_{a\in A}x^{\left\vert
a\right\vert }\right)  }_{=\overline{A}}\underbrace{\left(  \sum_{b\in
B}x^{\left\vert b\right\vert }\right)  }_{=\overline{B}}=\overline{A}%
\cdot\overline{B}.
\end{align*}
The proof of Proposition \ref{prop.gf-ws.prod} is thus complete.
\end{proof}

We can easily extend Definition \ref{def.gf-ws.prod} and Proposition
\ref{prop.gf-ws.prod} to Cartesian products of $k$ weighted sets. The weight
function on such a Cartesian product $A_{1}\times A_{2}\times\cdots\times
A_{k}$ is defined by%
\begin{equation}
\left\vert \left(  a_{1},a_{2},\ldots,a_{k}\right)  \right\vert =\left\vert
a_{1}\right\vert +\left\vert a_{2}\right\vert +\cdots+\left\vert
a_{k}\right\vert . \label{eq.gf-ws.prodk.weight}%
\end{equation}
As a particular case of such Cartesian products, we obtain Cartesian powers
when we multiply $k$ copies of the same weighted set:

\begin{definition}
Let $A$ be a weighted set. Then, $A^{k}$ (for $k\in\mathbb{N}$) means the
weighted set $\underbrace{A\times A\times\cdots\times A}_{k\text{ times}}$.
\end{definition}

The analogue of Proposition \ref{prop.gf-ws.prod} for Cartesian products of
$k$ weighted sets then yields the following:

\begin{proposition}
\label{prop.gf-ws.pow}Let $A$ be a finite-type weighted set. Let
$k\in\mathbb{N}$. Then, $A^{k}$ is finite-type, too, and satisfies
$\overline{A^{k}}=\overline{A}^{k}$.
\end{proposition}

Note that the $0$-th Cartesian power $A^{0}$ of a weighted set $A$ always
consists of a single element -- namely, the empty $0$-tuple $\left(
{}\right)  $, which has weight $0$.

If $A$ is a weighted set, then the infinite disjoint union $A^{0}+A^{1}%
+A^{2}+\cdots$ consists of all (finite) tuples of elements of $A$ (including
the $0$-tuple, the $1$-tuples, and so on). The weight of a tuple is the sum of
the weights of its entries (indeed, this is just what
(\ref{eq.gf-ws.prodk.weight}) says).

\subsubsection{Examples}

Now, let us use this theory to revisit some of the things we have already counted:

\begin{itemize}
\item Fix $k\in\mathbb{N}$, and let
\begin{align*}
C_{k}:=  &  \ \left\{  \text{compositions of length }k\right\} \\
=  &  \ \left\{  \left(  a_{1},a_{2},\ldots,a_{k}\right)  \ \mid\ a_{1}%
,a_{2},\ldots,a_{k}\text{ are positive integers}\right\} \\
=  &  \ \mathbb{P}^{k},\ \ \ \ \ \ \ \ \ \ \text{where }\mathbb{P}=\left\{
1,2,3,\ldots\right\}  .
\end{align*}
This set $C_{k}$ becomes a finite-type weighted set if we set $\left\vert
\left(  a_{1},a_{2},\ldots,a_{k}\right)  \right\vert =a_{1}+a_{2}+\cdots
+a_{k}$ for every $\left(  a_{1},a_{2},\ldots,a_{k}\right)  \in C_{k}$. What
is its weight generating function $\overline{C_{k}}$ ? We can turn
$\mathbb{P}$ itself into a weighted set, by defining the weight of a positive
integer $n$ by $\left\vert n\right\vert =n$. Then, $C_{k}=\mathbb{P}^{k}$ not
just as sets, but as weighted sets\footnote{\textit{Proof.} The weight of a
composition $\left(  a_{1},a_{2},\ldots,a_{k}\right)  \in C_{k}$ in $C_{k}$ is%
\[
\left\vert \left(  a_{1},a_{2},\ldots,a_{k}\right)  \right\vert =a_{1}%
+a_{2}+\cdots+a_{k}\ \ \ \ \ \ \ \ \ \ \left(  \text{by definition}\right)  .
\]
The weight of a composition $\left(  a_{1},a_{2},\ldots,a_{k}\right)
\in\mathbb{P}^{k}$ in $\mathbb{P}^{k}$ is%
\begin{align*}
\left\vert \left(  a_{1},a_{2},\ldots,a_{k}\right)  \right\vert  &
=\left\vert a_{1}\right\vert +\left\vert a_{2}\right\vert +\cdots+\left\vert
a_{k}\right\vert \ \ \ \ \ \ \ \ \ \ \left(  \text{by
(\ref{eq.gf-ws.prodk.weight})}\right) \\
&  =a_{1}+a_{2}+\cdots+a_{k}\ \ \ \ \ \ \ \ \ \ \left(  \text{since
}\left\vert n\right\vert =n\text{ for each }n\in\mathbb{P}\right)  .
\end{align*}
These two weights are clearly equal. Thus, the weight functions of $C_{k}$ and
$\mathbb{P}^{k}$ agree. Hence, $C_{k}=\mathbb{P}^{k}$ as weighted sets.}.
Hence,%
\begin{align*}
\overline{C_{k}}  &  =\overline{\mathbb{P}^{k}}=\overline{\mathbb{P}}%
^{k}\ \ \ \ \ \ \ \ \ \ \left(  \text{by Proposition \ref{prop.gf-ws.pow}%
}\right) \\
&  =\left(  x^{1}+x^{2}+x^{3}+\cdots\right)  ^{k}\ \ \ \ \ \ \ \ \ \ \left(
\text{since }\overline{\mathbb{P}}=x^{1}+x^{2}+x^{3}+\cdots\right) \\
&  =\left(  \dfrac{x}{1-x}\right)  ^{k}\ \ \ \ \ \ \ \ \ \ \left(  \text{since
}x^{1}+x^{2}+x^{3}+\cdots=\dfrac{x}{1-x}\right)  .
\end{align*}
This recovers the equality (\ref{pf.thm.fps.comps.num-comps-n-k.Ak=}).

\item Recall the notion of Dyck paths (as defined in Example 2 in Section
\ref{sec.gf.exas}), as well as the Catalan numbers $c_{0},c_{1},c_{2},\ldots$
(defined in the same place). Let
\[
D:=\left\{  \text{Dyck paths from }\left(  0,0\right)  \text{ to }\left(
2n,0\right)  \text{ for some }n\in\mathbb{N}\right\}  .
\]
This set $D$ becomes a finite-type weighted set if we set
\[
\left\vert p\right\vert =n\ \ \ \ \ \ \ \ \ \ \text{whenever }p\text{ is a
Dyck path from }\left(  0,0\right)  \text{ to }\left(  2n,0\right)  .
\]
The weight generating function $\overline{D}$ of this weighted set $D$ is%
\[
\overline{D}=\sum_{n\in\mathbb{N}}\underbrace{\left(  \text{\# of Dyck paths
from }\left(  0,0\right)  \text{ to }\left(  2n,0\right)  \right)
}_{\substack{=c_{n}\\\text{(by the definition of }c_{n}\text{)}}}x^{n}%
=\sum_{n\in\mathbb{N}}c_{n}x^{n}.
\]
This is the generating function we called $C\left(  x\right)  $ in Example 2
of Section \ref{sec.gf.exas}.

Recall that there is only one Dyck path from $\left(  0,0\right)  $ to
$\left(  0,0\right)  $, namely the trivial path. All the other Dyck paths in
$D$ are nontrivial. We let%
\begin{align*}
D_{\operatorname*{triv}}  &  :=\left\{  \text{trivial Dyck paths in
}D\right\}  \ \ \ \ \ \ \ \ \ \ \text{and}\\
D_{\operatorname*{non}}  &  :=\left\{  \text{nontrivial Dyck paths in
}D\right\}  .
\end{align*}
These two sets $D_{\operatorname*{triv}}$ and $D_{\operatorname*{non}}$ are
subsets of $D$, and thus are weighted sets themselves (we define their weight
functions by restricting the one of $D$). The set $D_{\operatorname*{triv}}$
consists of a single Dyck path, which has weight $0$; thus, its weight
generating function is%
\[
\overline{D_{\operatorname*{triv}}}=x^{0}=1.
\]


In Example 2 of Section \ref{sec.gf.exas}, we have seen that any nontrivial
Dyck path $\pi$ has the following structure:\footnote{The colors are referring
to the following picture:%
\[%
% [tikz figure removed]
%BeginExpansion
% [tikz figure removed]%
%EndExpansion
\]
}

\begin{itemize}
\item a NE-step,

\item followed by a (diagonally shifted) Dyck path (drawn in green),

\item followed by a SE-step,

\item followed by another (horizontally shifted) Dyck path (drawn in purple).
\end{itemize}

If we denote the green Dyck path by $\alpha$ and the purple Dyck path by
$\beta$, then we obtain a pair $\left(  \alpha,\beta\right)  \in D\times D$ of
two Dyck paths. Thus, each nontrivial Dyck path $\pi\in D_{\operatorname*{non}%
}$ gives rise to a pair $\left(  \alpha,\beta\right)  \in D\times D$ of two
Dyck paths. This yields a map%
\begin{align*}
D_{\operatorname*{non}} &  \rightarrow D\times D,\\
\pi &  \mapsto\left(  \alpha,\beta\right)  ,
\end{align*}
which is easily seen to be a bijection (since any pair $\left(  \alpha
,\beta\right)  \in D\times D$ can be assembled into a single nontrivial Dyck
path $\pi$ that starts with an NE-step, is followed by a shifted copy of
$\alpha$, then by a SE-step, then by a shifted copy of $\beta$).

Alas, this bijection is not an isomorphism of weighted sets, since it fails to
preserve the weight. Indeed, $\left\vert \left(  \alpha,\beta\right)
\right\vert =\left\vert \alpha\right\vert +\left\vert \beta\right\vert
=\left\vert \pi\right\vert -1\neq\left\vert \pi\right\vert $.

Fortunately, we can fix this rather easily. Define a weighted set%
\[
X:=\left\{  1\right\}  \ \ \ \ \ \ \ \ \ \ \text{with }\left\vert 1\right\vert
=1.
\]
This is a one-element set, so the only real difference between the weighted
sets $X\times D\times D$ and $D\times D$ is in the weights. Indeed, the sets
$D\times D$ and $X\times D\times D$ are in bijection (any pair $\left(
\alpha,\beta\right)  \in D\times D$ corresponds to the triple $\left(
1,\alpha,\beta\right)  \in X\times D\times D$), but the weights of
corresponding elements differ by $1$ (namely, $\left\vert \left(
1,\alpha,\beta\right)  \right\vert =\underbrace{\left\vert 1\right\vert }%
_{=1}+\underbrace{\left\vert \alpha\right\vert +\left\vert \beta\right\vert
}_{=\left\vert \left(  \alpha,\beta\right)  \right\vert }=1+\left\vert \left(
\alpha,\beta\right)  \right\vert $).

Thus, by replacing $D\times D$ by $X\times D\times D$, we can fix the degrees
in our above bijection. We thus obtain a bijection%
\begin{align*}
D_{\operatorname*{non}} &  \rightarrow X\times D\times D,\\
\pi &  \mapsto\left(  1,\alpha,\beta\right)  ,
\end{align*}
which does preserve the weight. This bijection is thus an isomorphism of
weighted sets. Hence, $D_{\operatorname*{non}}\cong X\times D\times D$.

Each Dyck path is either trivial or nontrivial. Hence,%
\[
D\cong D_{\operatorname*{triv}}+\underbrace{D_{\operatorname*{non}}}_{\cong
X\times D\times D}\cong D_{\operatorname*{triv}}+X\times D\times D,
\]
so that%
\begin{align*}
\overline{D} &  =\overline{D_{\operatorname*{triv}}+X\times D\times
D}\ \ \ \ \ \ \ \ \ \ \left(  \text{by Proposition \ref{prop.gf-ws.iso}%
}\right)  \\
&  =\underbrace{\overline{D_{\operatorname*{triv}}}}_{=1}%
+\underbrace{\overline{X}}_{=x}\cdot\underbrace{\overline{D}\cdot\overline{D}%
}_{=\overline{D}^{2}}\ \ \ \ \ \ \ \ \ \ \left(  \text{by Proposition
\ref{prop.gf-ws.djun} and Proposition \ref{prop.gf-ws.prod}}\right)  \\
&  =1+x\cdot\overline{D}^{2}.
\end{align*}
This is precisely the quadratic equation $C\left(  x\right)  =1+x\left(
C\left(  x\right)  \right)  ^{2}$ that we obtained in Section
\ref{sec.gf.exas}. But this time, we obtained it in a more conceptual way,
through a combinatorially defined isomorphism $D\cong D_{\operatorname*{triv}%
}+X\times D\times D$ rather than by ad-hoc manipulation of FPSs.
\end{itemize}

\subsubsection{\label{subsec.gf.weighted-set.domino}Domino tilings}

Let us now apply the theory of weight generating functions to something we
haven't already counted. Namely, we shall count the domino tilings of a
rectangle. Informally, these are defined as follows:

\begin{itemize}
\item For any $n,m\in\mathbb{N}$, we let $R_{n,m}$ be a rectangle with width
$n$ and height $m$.

\item A \emph{domino} means a rectangle that is an $R_{1,2}$ or an $R_{2,1}$.

\item A \emph{domino tiling} of a shape $S$ means a tiling of $S$ by dominos
(i.e., a set of dominos that cover $S$ and whose interiors don't intersect).
\end{itemize}

For example, here is a domino tiling of $R_{8,8}$:%
\[%
% [tikz figure removed]
%BeginExpansion
% [tikz figure removed]%
%EndExpansion
\]
(note that the colors are purely ornamental here: we are coloring a domino
pink if it lies horizontally and green if it stands vertically, for the sake
of convenience).

This sounds geometric, but actually is a combinatorial object hiding behind
geometric language. Our rectangles and dominos all align to a square grid.
Thus, rectangles can be modeled simply as finite sets of grid squares, and
dominos are unordered pairs of adjacent grid squares. Grid squares, in turn,
can be modeled as pairs $\left(  i,j\right)  $ of integers (corresponding to
the Cartesian coordinates of their centers). Thus, we redefine domino tilings
combinatorially as follows:

\begin{definition}
\label{def.domino.shapes-and-tilings}\textbf{(a)} A \emph{shape} means a
subset of $\mathbb{Z}^{2}$.

We draw each $\left(  i,j\right)  \in\mathbb{Z}^{2}$ as a unit square with
center at the point $\left(  i,j\right)  $ (in Cartesian coordinates); thus, a
shape can be drawn as a cluster of squares. \medskip

\textbf{(b)} For any $n,m\in\mathbb{N}$, the shape $R_{n,m}$ (called the
$n\times m$\emph{-rectangle}) is defined to be%
\[
\left\{  1,2,\ldots,n\right\}  \times\left\{  1,2,\ldots,m\right\}  =\left\{
\left(  i,j\right)  \in\mathbb{Z}^{2}\ \mid\ 1\leq i\leq n\text{ and }1\leq
j\leq m\right\}  .
\]


\textbf{(c)} A \emph{domino} means a size-$2$ shape of the form%
\begin{align*}
&  \left\{  \left(  i,j\right)  ,\ \left(  i+1,j\right)  \right\}  \text{ (a
\textquotedblleft\emph{horizontal domino}\textquotedblright)}%
\ \ \ \ \ \ \ \ \ \ \text{or}\\
&  \left\{  \left(  i,j\right)  ,\ \left(  i,j+1\right)  \right\}  \text{ (a
\textquotedblleft\emph{vertical domino}\textquotedblright)}%
\end{align*}
for some $\left(  i,j\right)  \in\mathbb{Z}^{2}$. \medskip

\textbf{(d)} A \emph{domino tiling} of a shape $S$ is a set partition of $S$
into dominos (i.e., a set of disjoint dominos whose union is $S$). \medskip

\textbf{(e)} For any $n,m\in\mathbb{N}$, let $d_{n,m}$ be the \# of domino
tilings of $R_{n,m}$.
\end{definition}

For example, the domino tiling $\raisebox{-3.5ex}{
% [tikz figure removed]
}$ of the rectangle $R_{3,2}$ is the set partition%
\[
\left\{  \left\{  \left(  1,1\right)  ,\ \left(  1,2\right)  \right\}
,\ \left\{  \left(  2,1\right)  ,\ \left(  3,1\right)  \right\}  ,\ \left\{
\left(  2,2\right)  ,\ \left(  3,2\right)  \right\}  \right\}
\]
of $R_{3,2}$ (here, $\left\{  \left(  1,1\right)  ,\ \left(  1,2\right)
\right\}  $ is the vertical domino, whereas $\left\{  \left(  2,1\right)
,\ \left(  3,1\right)  \right\}  $ is the bottom horizontal domino, and
$\left\{  \left(  2,2\right)  ,\ \left(  3,2\right)  \right\}  $ is the top
horizontal domino). \medskip

Can we compute $d_{n,m}$ ?

The case $m=1$ is a bit too simple (do it!), so let us start with the case
$m=2$. Here are the $d_{n,2}$ for $n\in\left\{  0,1,\ldots,4\right\}  $:%
\[%
\begin{tabular}
[c]{|c|c|l|}\hline
$n$ & $d_{n,2}$ & domino tilings of $R_{n,2}$\\\hline\hline
$0$ & $d_{0,2}=1$ &
$\vphantom{\dfrac{\int}{\int}}% [tikz figure removed]$%
\\\hline
$1$ & $d_{1,2}=1$ &
$\vphantom{\dfrac{\int}{\int}}% [tikz figure removed]$%
\\\hline
$2$ & $d_{2,2}=2$ &
$\vphantom{\dfrac{\int}{\int}}% [tikz figure removed]\ \ \ ,\ \ \ % [tikz figure removed]$%
\\\hline
$3$ & $d_{3,2}=3$ &
$\vphantom{\dfrac{\int}{\int}}% [tikz figure removed]\ \ \ ,\ \ \ % [tikz figure removed]\ \ \ ,\ \ \ % [tikz figure removed]$%
\\\hline
$4$ & $d_{4,2}=5$ &
$\vphantom{\dfrac{\int}{\int}}% [tikz figure removed]\ \ \ ,\ \ \ % [tikz figure removed]\ \ \ ,\ \ \ % [tikz figure removed]\ \ \ ,$%
\\
&  &
$\vphantom{\dfrac{\int}{\int}}% [tikz figure removed]\ \ \ ,\ \ \ % [tikz figure removed]$%
\\\hline
\end{tabular}
\ \ \ \ \ \ \ \ \ .
\]


Let us try to compute $d_{n,2}$ in general.

A \emph{height-}$2$ \emph{rectangle} shall mean a rectangle of the form
$R_{n,2}$ with $n\in\mathbb{N}$. Let us define the weighted set%
\begin{align*}
D:=  &  \left\{  \text{domino tilings of height-}2\text{ rectangles}\right\}
\\
=  &  \left\{  \text{domino tilings of }R_{n,2}\text{ with }n\in
\mathbb{N}\right\}  .
\end{align*}
We define the weight of a tiling $T$ of $R_{n,2}$ to be $\left\vert
T\right\vert :=n$ (that is, the width of the rectangle tiled by this tiling).

Thus, $D$ is a finite-type weighted set, with weight generating function%
\[
\overline{D}=\sum_{n\in\mathbb{N}}d_{n,2}x^{n}.
\]


So we want to compute $\overline{D}$. Let us define a new weighted set that
will help us at that.

Namely, we say that a \emph{fault} of a domino tiling $T$ is a vertical line
$\ell$ such that

\begin{itemize}
\item each domino of $T$ lies either left of $\ell$ or right of $\ell$ (but
does not straddle $\ell$), and

\item there is at least one domino of $T$ that lies left of $\ell$, and at
least one domino of $T$ that lies right of $\ell$.
\end{itemize}

For example, here are two domino tilings:%
\[%
% [tikz figure removed]
%BeginExpansion
% [tikz figure removed]%
%EndExpansion
\ \ \ \ \ \ \ \ \ \ \
% [tikz figure removed]
%BeginExpansion
% [tikz figure removed]%
%EndExpansion
\ \ \ \ \ \ .
\]
The tiling on the left has a fault (namely, the vertical line separating the
$2$nd from the $3$rd column), but the tiling on the right has none (a fault
must be a vertical line by definition; a horizontal line doesn't count).

A domino tiling will be called \emph{faultfree} if it is nonempty and has no
fault. Thus, the tiling on the right (in the above example) is faultfree.

We now observe a crucial (but trivial) lemma:

\begin{lemma}
\label{lem.gf.weighted-set.domino.fd}Any domino tiling of a height-$2$
rectangle can be decomposed uniquely into a tuple of faultfree tilings of
(usually smaller) height-$2$ rectangles, by cutting it along its faults. For
example:%
\begin{align*}
&  \raisebox{-3.5ex}{
% [tikz figure removed]
}\ \text{ decomposes as}\\
&  \left(  \ \ \raisebox{-3.5ex}{
% [tikz figure removed]
}\ \ \ \ \ ,\ \ \raisebox{-3.5ex}{
% [tikz figure removed]
}\ \ \ \ \ ,\ \ \raisebox{-3.5ex}{
% [tikz figure removed]
}\ \ \ \ \ ,\ \ \raisebox{-3.5ex}{
% [tikz figure removed]
}\ \ \right)  .
\end{align*}
(Note that if the original tiling was faultfree, then it will decompose into a
$1$-tuple. If the original tiling was empty, then it will decompose into a $0$-tuple.)

Moreover, the sum of the weights of the faultfree tilings in the tuple is the
weight of the original tiling. (In other words, if a tiling $T$ decomposes
into the tuple $\left(  T_{1},T_{2},\ldots,T_{k}\right)  $, then $\left\vert
T\right\vert =\left\vert T_{1}\right\vert +\left\vert T_{2}\right\vert
+\cdots+\left\vert T_{k}\right\vert $.)
\end{lemma}

Thus, if we define a new weighted set%
\[
F:=\left\{  \text{\textbf{faultfree} domino tilings of }R_{n,2}\text{ with
}n\in\mathbb{N}\right\}
\]
(with the same weights as in $D$), then we obtain an isomorphism (i.e.,
weight-preserving bijection)%
\[
D\rightarrow\underbrace{F^{0}+F^{1}+F^{2}+F^{3}+\cdots}_{\substack{\text{an
infinite disjoint union}\\\text{of weighted sets}}}
\]
that sends each tiling to the tuple it decomposes into. Hence,%
\[
D\cong F^{0}+F^{1}+F^{2}+F^{3}+\cdots,
\]
and therefore%
\begin{align*}
\overline{D}  &  =\overline{F^{0}+F^{1}+F^{2}+F^{3}+\cdots}=\overline{F^{0}%
}+\overline{F^{1}}+\overline{F^{2}}+\overline{F^{3}}+\cdots\\
&  \ \ \ \ \ \ \ \ \ \ \ \ \ \ \ \ \ \ \ \ \left(  \text{by the infinite
analogue of Proposition \ref{prop.gf-ws.djun}}\right) \\
&  =\overline{F}^{0}+\overline{F}^{1}+\overline{F}^{2}+\overline{F}^{3}%
+\cdots\ \ \ \ \ \ \ \ \ \ \left(  \text{by Proposition \ref{prop.gf-ws.pow}%
}\right) \\
&  =\dfrac{1}{1-\overline{F}}\ \ \ \ \ \ \ \ \ \ \left(  \text{by
(\ref{eq.sec.gf.exas.1.1/1-x}), with }\overline{F}\text{ substituted for
}x\right)  .
\end{align*}
Thus, if we can compute $\overline{F}$, then we can compute $\overline{D}$.

In order to compute $\overline{F}$, let us see how a faultfree domino tiling
of a height-$2$ rectangle looks like. Here are two such tilings:%
\[
\raisebox{-3.5ex}{
% [tikz figure removed]
}\ \ \ \ \ \ \ \ \ \ \text{and}\ \ \ \ \ \ \ \ \ \ \raisebox{-3.5ex}{
% [tikz figure removed]
}\ \ \ \ \ .
\]


I claim that these two tilings are the \textbf{only} faultfree tilings of
height-$2$ rectangles. Indeed, consider any faultfree tiling of a height-$2$
rectangle. In this tiling, look at the domino that covers the box $\left(
1,1\right)  $. If it is a vertical domino, then this vertical domino must
constitute the entire tiling, since otherwise there would be a fault to its
right. If it is a horizontal domino, then there must be a second horizontal
domino stacked atop it, and these two dominos must then constitute the entire
tiling, since otherwise there would be a fault to their right. This leads to
the two options we just named.

Thus, the weighted set $F$ consists of just the two tilings shown above: one
tiling of weight $1$ and one tiling of weight $2$. Hence, its weight
generating function is $\overline{F}=x+x^{2}$. So%
\[
\overline{D}=\dfrac{1}{1-\overline{F}}=\dfrac{1}{1-\left(  x+x^{2}\right)
}=\dfrac{1}{1-x-x^{2}}=f_{1}+f_{2}x+f_{3}x^{2}+f_{4}x^{3}+\cdots,
\]
where $\left(  f_{0},f_{1},f_{2},\ldots\right)  $ is the Fibonacci sequence.
Thus, comparing coefficients, we find%
\[
d_{n,2}=f_{n+1}\ \ \ \ \ \ \ \ \ \ \text{for each }n\in\mathbb{N}.
\]


There are, of course, more elementary proofs of this (see \cite[Proposition
1.1.11]{19fco}).

\begin{remark}
Here is an outline of an alternative proof of $\overline{D}=\dfrac
{1}{1-\overline{F}}$:

Any tiling in $D$ is either empty, or can be uniquely split into a pair of a
faultfree tiling and an arbitrary tiling (just split it along its leftmost
fault, or along the right end if there is no fault). Thus,%
\[
D\cong\left\{  0\right\}  +F\times D,
\]
where the set $\left\{  0\right\}  $ here is viewed as a weighted set with
$\left\vert 0\right\vert =0$. Hence,%
\[
\overline{D}=\overline{\left\{  0\right\}  +F\times D}=\overline{\left\{
0\right\}  }+\overline{F}\cdot\overline{D}=1+\overline{F}\cdot\overline{D}.
\]
Solving this for $\overline{D}$, we find $\overline{D}=\dfrac{1}%
{1-\overline{F}}.$
\end{remark}

Now, let us try to solve the analogous problem for height-$3$ rectangles.
Forget about the $D$ and $F$ we defined above. Instead, define a new weighted
set%
\begin{align*}
D:=  &  \left\{  \text{domino tilings of height-}3\text{ rectangles}\right\}
\\
=  &  \left\{  \text{domino tilings of }R_{n,3}\text{ with }n\in
\mathbb{N}\right\}  .
\end{align*}
The weight of a tiling $T$ of $R_{n,3}$ is defined as before (i.e., it is
$\left\vert T\right\vert =n$). Thus, $D$ is a finite-type weighted set, with
generating function%
\[
\overline{D}=\sum_{n\in\mathbb{N}}d_{n,3}x^{n}.
\]
We want to compute this $\overline{D}$.

Set%
\[
F:=\left\{  \text{\textbf{faultfree} domino tilings of }R_{n,3}\text{ with
}n\in\mathbb{N}\right\}  .
\]
Again, we can show $\overline{D}=\dfrac{1}{1-\overline{F}}$ (as we did above
in the case of $R_{n,2}$). We thus need to compute $\overline{F}$.

How does a faultfree domino tiling of a height-$3$ rectangle look like? Let us
classify such tilings according to the kind of dominos that occupy the first
column of the tiling. The proof of the following classification can be found
in the Appendix (Proposition \ref{prop.gf.weighted-set.domino.Rn3.ABC} in
Section \ref{sec.details.domino}):

\begin{itemize}
\item The faultfree domino tilings of a height-$3$ rectangle that contain a
vertical domino in the \textbf{top} two squares of the first column are%
\[%
% [tikz figure removed]
%BeginExpansion
% [tikz figure removed]
\qquad% [tikz figure removed]
\qquad% [tikz figure removed]%
%EndExpansion
\]
(this is an infinite sequence of tilings, each obtained from the previous by
inserting two columns in the middle by a fairly self-explanatory procedure).
The weights of these tilings are $2,4,6,\ldots$, so their total contribution
to the weight generating function $\overline{F}$ of $F$ is $x^{2}+x^{4}%
+x^{6}+\cdots$.

\item The faultfree domino tilings of a height-$3$ rectangle that contain a
vertical domino in the \textbf{bottom} two squares of the first column are%
\[%
% [tikz figure removed]
%BeginExpansion
% [tikz figure removed]
\qquad% [tikz figure removed]
\qquad% [tikz figure removed]%
%EndExpansion
\]
(these are the top-down mirror images of the tilings from the previous bullet
point). The weights of these tilings are $2,4,6,\ldots$, so their total
contribution to the weight generating function $\overline{F}$ of $F$ is
$x^{2}+x^{4}+x^{6}+\cdots$.

\item The faultfree domino tilings of a height-$3$ rectangle that contain
\textbf{no} vertical domino in the first column are%
\[%
% [tikz figure removed]
%BeginExpansion
% [tikz figure removed]%
%EndExpansion
\]
(yes, there is only one such tiling). The weight of this tiling is $2$, so its
total contribution to the weight generating function $\overline{F}$ of $F$ is
$x^{2}$.
\end{itemize}

This classification of faultfree domino tilings entails%
\begin{align*}
\overline{F}  &  =\left(  x^{2}+x^{4}+x^{6}+\cdots\right)  +\left(
x^{2}+x^{4}+x^{6}+\cdots\right)  +x^{2}\\
&  =x^{2}\cdot\dfrac{1}{1-x^{2}}+x^{2}\cdot\dfrac{1}{1-x^{2}}+x^{2}%
\ \ \ \ \ \ \ \ \ \ \left(  \text{since }x^{2}+x^{4}+x^{6}+\cdots=x^{2}%
\cdot\dfrac{1}{1-x^{2}}\right) \\
&  =\dfrac{3x^{2}-x^{4}}{1-x^{2}}.
\end{align*}
Thus,%
\begin{align*}
\overline{D}  &  =\dfrac{1}{1-\overline{F}}=\dfrac{1}{1-\dfrac{3x^{2}-x^{4}%
}{1-x^{2}}}=\dfrac{1-x^{2}}{1-4x^{2}+x^{4}}\\
&  =1+3x^{2}+11x^{4}+41x^{6}+153x^{8}+\cdots.
\end{align*}


You will notice that only even powers of $x$ appear in this FPS. In other
words,%
\[
d_{n,3}=0\text{ when }n\text{ is odd.}%
\]
This is not surprising, because if $n$ is odd, then the rectangle $R_{n,3}$
has an odd \# of squares, and thus cannot be tiled by dominos.

But we can also compute $d_{n,3}$ for even $n$. Indeed, using the same method
(partial fractions) that we used for the Fibonacci sequence in Section
\ref{sec.gf.exas}, we can expand $\dfrac{1-x^{2}}{1-4x^{2}+x^{4}}$ as a sum of
geometric series:%
\[
\dfrac{1-x^{2}}{1-4x^{2}+x^{4}}=\dfrac{3+\sqrt{3}}{6}\cdot\dfrac{1}{1-\left(
2+\sqrt{3}\right)  x^{2}}+\dfrac{3-\sqrt{3}}{6}\cdot\dfrac{1}{1-\left(
2-\sqrt{3}\right)  x^{2}}.
\]
Thus, we find%
\[
d_{n,3}=\dfrac{3+\sqrt{3}}{6}\left(  2+\sqrt{3}\right)  ^{n/2}+\dfrac
{3-\sqrt{3}}{6}\left(  2-\sqrt{3}\right)  ^{n/2}\ \ \ \ \ \ \ \ \ \ \text{for
any even }n.
\]


Now, what about computing $d_{n,m}$ in general? The above reasoning leading up
to $\overline{D}=\dfrac{1}{1-\overline{F}}$ can be applied for any
$m\in\mathbb{N}$, but describing $F$ becomes harder and harder as $m$ grows
larger. The generating function $\overline{D}$ is still a quotient of two
polynomials for any $m$ (see, e.g., \cite{KlaPol79}), but this requires more
insight to prove. For $m\geq6$, it appears that there is no formula for
$d_{n,m}$ that requires only quadratic irrationalities.

It is worth mentioning a different formula for $d_{n,m}$, found by Kasteleyn
in 1961 (motivated by a theoretical physics model):

\begin{theorem}
[Kasteleyn's formula]Let $n,m\in\mathbb{N}$ be such that $m$ is even and
$n\geq1$. Then,%
\[
d_{n,m}=2^{mn/2}\ \prod_{j=1}^{m/2}\ \ \prod_{k=1}^{n}\sqrt{\left(  \cos
\dfrac{j\pi}{m+1}\right)  ^{2}+\left(  \cos\dfrac{k\pi}{n+1}\right)  ^{2}}.
\]

\end{theorem}

See \cite[Theorem 12.85]{Loehr-BC} or \cite{Stucky15} for proofs of this
rather surprising formula (which, alas, require some more advanced methods
than those introduced in this text). Note that it can indeed be used for exact
computation of $d_{n,m}$ (as there are algorithms for exact manipulation of
\textquotedblleft trigonometric irrationals\textquotedblright\footnote{The
technical term is \textquotedblleft elements of cyclotomic
fields\textquotedblright.}\ such as $\cos\dfrac{j\pi}{m+1}$); for example, it
yields $d_{8,8}=12\ 988\ 816$.
